\chapter{Conclusion and Future Work}
\label{ch:conclusion}

\section{Summary}
\label{sec:c-summary}

This thesis asked whether a mechanistic-interpretability readout of an unlearned language model's internal
computation can yield an evaluation metric that measures knowledge rather than surface behaviour, and
whether such a metric can be shown faithful. Working on the Hubble benchmark, whose controlled fact
insertions provide a never-knew and a knows reference for every fact, we integrated a corrected,
boundary-safe logit lens into unlearning evaluation and derived from its rank-by-layer trajectory a
dual-anchor distance, $D_\text{mech}$.

The results answer the four research questions. Output and lexical metrics collapse the four unlearning
methods onto a single ``forgotten'' verdict (RQ1); the internal readout resolves that verdict into four
distinct mechanisms and shows that none of them removes the fact's internal trace under the readouts
applied, so on this benchmark unlearning suppresses at the output rather than deleting the representation
(RQ2). Under the faithful analysis based on OpenUnlearning, on Hubble-8B $D_\text{mech}$ is the only metric,
mechanistic or lexical, that is both strongly knowledge-discriminative and substantially paraphrase-stable,
with a composite faithfulness score far exceeding any lexical metric; the win comes from the
relative anchor-difference construction, not merely from reading internals, since a na\"ive internal metric
fails paraphrase stability just as the lexical metrics do, and the ranking survives every aggregation we
tried. We replicate the meta-evaluation on OpenUnlearning TOFU at 1B, 3B, and 8B, where $D_\text{mech}$ is
again faithful and the best internal metric at every scale, though a few robust lexical aggregates become
competitive (RQ3). Finally, comparing the three lenses shows why the metric is built on the parameter-free
logit lens rather than the calibrated tuned lens, which ``skips ahead'' and hides suppression, or the
interventionally-grounded Jacobian lens, which is a valuable confirmatory readout but a less suitable metric
basis; and the whole picture is strongly scale-dependent, with several methods internalising unlearning at
8B but not at 1B (RQ4).

The broader contribution is a proof of concept that unlearning evaluation can, and arguably should, read the
model's internal state. As unlearning becomes a compliance instrument for privacy, copyright, and safety,
the gap between looking forgotten and being forgotten is exactly the gap a faithful metric must close, and a
metric that a paraphrase can defeat is not measuring knowledge.

\section{Future Work}
\label{sec:c-future}

Several directions would strengthen and extend these results.

\paragraph{Hardening the faithfulness claim.}
The most immediate improvement is to replace the single fixed paraphrase and wrong-target with samples of
many machine-generated rewordings and distractor styles, so that stability is measured across paraphrase
\emph{styles} as well as across facts. The cross-benchmark replication on TOFU at 1B, 3B, and 8B
(Section~\ref{sec:r-crossbench}) already extends the study beyond Hubble; further attributes beyond
birthdate and additional model families would harden the external-validity claim further.

\paragraph{From observation to causation.}
The thesis establishes suppression-versus-erasure observationally. Causal-intervention methods, such as
activation patching or steering of the suppressed representation, would turn ``the fact is still internally
present'' into a causal statement, by showing that restoring the mid-stack representation restores the
output. This would also connect the evaluation work back to the intervention-side interpretability
literature.

\paragraph{A stronger interventional lens.}
Our Jacobian lens averages over a modest reference set. Implementing the full averaged-transport pipeline of
\citet{anthropic2026workspace}, with its larger reference corpus and normalisation, may sharpen this
readout and, in particular, recover the knows-versus-never-knew separation at 1B that our lightweight
version misses.

\paragraph{Generality across benchmarks and models.}
Finally, the approach should be tested beyond Hubble. Applying the dual-anchor construction to TOFU, MUSE,
and WMDP would require constructing knows and never-knew references there, and would establish whether the
faithfulness advantage of an internal metric is a property of Hubble's controlled insertions or a general
feature of unlearning evaluation. Integrating $D_\text{mech}$, or a reference-free approximation of it, into
a standard toolkit such as OpenUnlearning \citep{dorna2025openunlearning} would make internal faithfulness a
routine part of how unlearning is reported.
